\documentclass[10pt,a4paper]{amsart}
\usepackage[margin=19mm]{geometry}
\usepackage{amsmath,amssymb,mathtools}
\usepackage[scr=rsfs,cal=euler]{mathalfa}
\usepackage{xfrac}
\usepackage[unicode,hidelinks,bookmarks=false]{hyperref}
\newcommand{\ie}{i.e.\ }
\newcommand{\cf}{cf.\ }
\newcommand{\resp}{respectively\ }
\newcommand{\Subsection}{\subsection}
\providecommand{\nobreakdash}{\nobreak\hbox{-}}
\setlength{\emergencystretch}{2em}
\multlinegap0pt
\usepackage{bm}
\usepackage{bbm}
\usepackage{dsfont}
\usepackage{xspace}
\usepackage{cancel}
\usepackage{mathtools}
\usepackage{cleveref}
\usepackage{amsthm}
\makeatletter
\@ifundefined{lem}{\newtheorem{lem}{Lem}}{}
\@ifundefined{prop}{\newtheorem{prop}{Prop}}{}
\makeatother
\title[Order-Preserving outer automorphisms of free and surface gro]{Order-Preserving outer automorphisms of free and surface groups}
\author{Jonathan Johnson, Khanh Le}
\date{Source: arXiv:2501.18847v2 (CC BY 4.0)}
\begin{document}
\maketitle

\noindent\emph{Source and licence.} Order-Preserving outer automorphisms of free and surface groups. Version arXiv:2501.18847v2 (CC BY 4.0), distributed under the Creative Commons Attribution 4.0 International licence, as recorded in the arXiv licence metadata for this exact version and shown on the abstract page https://arxiv.org/abs/2501.18847v2. Full rights notice: licenses/arxiv-2501.18847-CC-BY-4.0.txt.\par\smallskip

\noindent\textit{Editorial scope.} Counted unit: the Prop whose source label is \texttt{prop:OrderingOnTensorPowerOfPuiseuxField} in the article (source0.tex lines 678--681, source label \texttt{prop:OrderingOnTensorPowerOfPuiseuxField}), delivered together with the article's own complete original proof of it (source0.tex lines 683--766, one \texttt{proof} environment of 84 lines). No mathematics is written, repaired or re-derived: statement and proof are the article's own text line for line. Required context retained uncounted: none. Every definition, display and label of those lines is retained unchanged. Omitted from this package and not redistributed: 1--677, 682--682, 767--1313. Nothing inside a retained excerpt is omitted or altered: every retained body is delivered exactly as the article's lines are written, so no omission receipt of any kind accompanies this package. Citations: the retained excerpts carry no citation command at all, so no bibliography entry of the article is transcribed and no reference is redistributed. Reference adaptation: none was needed and none was made; every reference inside the delivered unit resolves inside it, so no unresolved reference or citation is produced. Printed slips kept literal: no printed word, symbol, hypothesis or number of the counted statement or of its proof was altered. Layout and class: the article's own class is \texttt{amsart} with packages including float, graphicx, JohnsonLe, todonotes, biblatex, lineno; the wrapper is the lane's pinned \texttt{amsart} editorial wrapper at 10pt on A4, which restates the theorem-like environments the retained text needs and adds the article's own macro definitions that the retained text needs (none), with extra wrapper packages (bm, bbm, dsfont, xspace, cancel, mathtools, cleveref). The delivered numbering is the unit's own. Assets: no retained span contains a figure environment, a graphics-inclusion command, a PDF-page inclusion, a table or a TikZ picture; no image file is copied into this package, the package declares no asset dependency and no image file is redistributed with it. Licence: version arXiv:2501.18847v2 of this article is distributed under the Creative Commons Attribution 4.0 International licence, as recorded in the arXiv licence metadata for this exact version and shown on the abstract page https://arxiv.org/abs/2501.18847v2. A full rights notice is delivered at licenses/arxiv-2501.18847-CC-BY-4.0.txt. Characters: every retained excerpt is pure ASCII, so the delivered engine, XeLaTeX with its default Latin Modern text font, reports no missing character.
\begin{prop}
    \label{prop:OrderingOnTensorPowerOfPuiseuxField}
    Let $m \in \mathbb{N}$. There exists an ordering $Q_m$ on $\mathbb{E}_m$ such that $u_1\otimes \dots \otimes u_m \in  Q_m$ for any $ \{u_i\}\subset Q$.
\end{prop}

\begin{proof}[Proof of \cref{prop:OrderingOnTensorPowerOfPuiseuxField}]
    When $m =1$, we have $\mathbb{E}_1 = \mathbb{E}$ and can simply take $Q_1 = Q$ as the ordering on $\mathbb{E}$. Fix a positive integer $m\geq 2$. We first define a lexicographic ordering $\prec$ on the following set of simple tensors
    \[
    S_m = \{t^{q_1} \otimes \dots \otimes t^{q_m} \mid q_i \in \mathbb{Q}\}.
    \]
    In particular, we write  
    \[
    t^{q_1} \otimes \dots \otimes t^{q_m} \prec t^{q'_1} \otimes \dots \otimes t^{q'_m} \Leftrightarrow q_{j} < q'_j \text{ in } \mathbb{Q} \text{ and } q_i = q'_i \ \forall \ 1 \leq i < j. 
    \]
    \begin{lem}
        The set $S_m$ equipped with $\prec$ is an ordered multiplicative abelian group.
    \end{lem}
    \begin{proof}
        Indeed, $S_m$ is a multiplicative abelian group with the group operation given by 
        \[
        (t^{q_1} \otimes \dots \otimes t^{q_m})(t^{q'_1} \otimes \dots \otimes t^{q'_m}) = (t^{q_1+q'_1} \otimes \dots \otimes t^{q_m+q'_m}).
        \]
        To see that $\prec$ defines a total bi-invariant ordering on $S_m$, we consider
        \begin{align*}
    P_{\prec} := \{s\in S_m \mid 1_{S_m} \prec s\} = \{t^{q_1} \otimes \dots \otimes t^{q_m} \mid \exists 1\leq j \leq m, q_j >0 \text{ and } \forall 1 \leq i <j, q_i = 0 \}.    
    \end{align*}
    This set is closed under multiplication and partitions $S_m = P_{\prec} \sqcup \{1_{S_m}\}\sqcup P_{\prec}^{-1}$. Therefore, $\prec$ defines a total bi-invariant ordering on $S_m$ when viewed as a multiplicative abelian group. 
    \end{proof}
    
    Now, we will use the ordering $\prec$ on $S_m$ to define an ordering $Q_m$ for $\mathbb{E}_m$. First, we make an observation about simple tensors in $\mathbb{E}_m$.

    \begin{lem}
        Let $u = u_1 \otimes \dots \otimes u_m \in \mathbb{E}_m$ be a simple tensor where
        \[
        u_i =\sum\limits_{j=k_i}^{\infty} c_{ij}t^{j/d_i} \in \mathbb{E}
        \]
        Then $u$ is a formal $\mathbb{R}$-linear combination of elements in $S_m$. Furthermore, if $u\neq 0$, then the set of all elements in $S_m$ with non-zero coefficient in $u$ is well-ordered with respect to $\prec$. 
    \end{lem}

    \begin{proof}
        Let $d$ be the l.c.m of $d_1,\dots,d_m$. We can rewrite each $u_i$ as 
        \[
        u_i = \sum\limits_{j=k_i}^{\infty} c_{ij}t^{j/d_i} = \sum\limits_{\ell_i=k'_i}^{\infty} b_{i\ell_i}t^{\ell_i/d}
        \]
        where $\ell_i = jd/d_i$ and $b_{i\ell_i} = c_{ij}$. Now we can write 
        \begin{align*}
            u 
            &= \left(\sum\limits_{\ell_1=k'_1}^{\infty} b_{1\ell_1}t^{\ell_1/d}\right) \otimes \dots \otimes \left(\sum\limits_{\ell_m=k'_m}^{\infty} b_{m\ell_m}t^{\ell_m/d}\right) \\
            &= \sum\limits_{\ell_1=k'_1}^{\infty}\dots\sum\limits_{\ell_m=k'_m}^{\infty}\left(b_{1\ell_1}\dots b_{m\ell_m}\right)\left(t^{\ell_1/d} \otimes \dots \otimes t^{\ell_m/d}\right).
        \end{align*}
         The calculation shows that any simple tensor is a formal $\mathbb{R}$-linear combination of elements in $S_m$. The set of all elements in $S_m$ with non-zero coefficient in $u$ is contained in 
        \[
        \{ t^{\ell_1/d} \otimes \dots \otimes t^{\ell_m/d} \mid \ell_i \geq k'_i \text{ and } \ell_i \in \mathbb{Z}\}.
        \]       
        This set is order-isomorphic to $\mathbb{N}^m$ with the lexicographic order and, therefore, is well-ordered. In particular, $t^{k'_1/d} \otimes \dots \otimes t^{k'_m/d}$ is the smallest element in $S_m$ with respect to $\prec$ that has a non-zero coefficient in $u$. 
    \end{proof}

    In view of the previous lemma, we can assume that any simple tensor in $\mathbb{E}_m$ has the form 
    \[
    u = \sum\limits_{j_1 \geq k_1,\dots,j_m \geq k_m}^\infty c_{j_1,\dots,j_m} (t^{j_1/d} \otimes \dots \otimes t^{j_m/d})
    \]
    where $c_{k_1,\dots,k_m} \neq 0$. In particular, the formal sum runs over the set 
    \[
        \{ t^{j_1/d} \otimes \dots \otimes t^{j_m/d} \mid j_i \geq k_i \text{ and } j_i \in \mathbb{Z}\}
    \]
    which is well-ordered with respect to $\prec$. 
    \begin{lem}
    \label{lem:WellDefinedLowestTerm}
        Let $f \in \mathbb{E}_m$. Then $f$ is a formal $\mathbb{R}$-linear combination of elements in $S_m$. Furthermore, if $f\neq 0$, then the set of all elements in $S_m$ with non-zero coefficient in $f$ is well-ordered with respect to $\prec$. Consequently, for any $f \neq 0$, the coefficient of the lowest term, $\mathfrak{c}_m(f) \in \mathbb{R}$, of $f$ with respect to $\prec$ is well-defined. 
    \end{lem}

    \begin{proof}
        The proof of this is similar to the previous lemma. The point is that the set of denominators is still finite since $f$ is a finite $\mathbb{R}$-linear combination of simple tensors. So we can rewrite everything in terms of a single common denominator. In particular, every element of $\mathbb{E}_m$ can be expressed as a formal linear combination of terms in the set 
        \[
        \{ t^{j_1/d} \otimes \dots \otimes t^{j_m/d} \mid j_i \geq k_i \text{ and } j_i \in \mathbb{Z}\}
        \] 
        which is well-ordered. 
    \end{proof}

    Now we are in position to define the ordering $Q_m$ on $\mathbb{E}_m$. Given a nonzero element $f\in \mathbb{E}_m$, we express $f$ as a formal sum of terms in $S_m$ with a single denominator. By \cref{lem:WellDefinedLowestTerm}, this formal sum has a well-defined lowest term among all nonzero terms with respect to the ordering $\prec $ on $S_m$. We let $\mathfrak{c}_m(f) \in \mathbb{R}$ be the coefficient of the lowest nonzero term of $f$ with respect to the ordering $\prec$ of $S_m$. In addition, $\mathfrak{c}_m(0)= 0$. We define
    \[
    Q_m := \{f \in \mathbb{E}_m \mid \mathfrak{c}_m(f) > 0 \}. 
    \]
    Since every nonzero element of $\mathbb{E}_m$ has a well-defined lowest term with respect to $\prec$, we have 
    \[\mathbb{E}_m = Q_m \sqcup\{0\} \sqcup -Q_m.\] 
    Now, let $f,g \in Q_m$. Since the ordering $\prec$ on $S_m$ is bi-invariant, the lowest term of $fg$ is the product of the lowest terms of $f$ and $g$. This implies that $\frak{c}_m(fg) = \frak{c}_m(f)\frak{c}_m(g) >0$ and that $Q_m Q_m \subseteq Q_m$. The lowest term of $f+g$ is either the lowest term of $f$, $g$ or the sum of both since $\mathfrak{c}_m(f), \mathfrak{c}_m(g)$ are both positive. Therefore, $\frak{c}_m(f+g)$ is either $\frak{c}_m(f)$, $\frak{c}_m(g)$, or $\frak{c}_m(f)+\frak{c}_m(g)$. In any case, $\frak{c}_m(f+g) >0$ which implies that $Q_m + Q_m \subset Q_m$. Hence, $Q_m$ defines an ordering on $\mathbb{E}_m$. 

    For the claim that $u_1\otimes \dots \otimes u_m \in  Q_m$, for any $ \{u_i\}\subset Q$, we consider $u_i \in Q \subset \mathbb{E}$ for $1\leq i \leq m$. Note that $\mathfrak{c}_m(u_1\otimes \dots \otimes u_m) = \frak{c}(u_1)\dots\frak{c}(u_m) >0$ which implies that $u_1\otimes \dots \otimes u_m \in Q_m$. This completes the proof of \cref{prop:OrderingOnTensorPowerOfPuiseuxField}.
\end{proof}

\end{document}
